
\noindent Durant mes années de recherche scientifique, j'ai principalement cherché à
résoudre deux problématiques très récurrentes du calcul haute performance, à
savoir l'accélération des calculs d'algèbre linéaire et la gestion de la
mémoire des données. Sans l'optimisation de ces deux aspects, la
résolution d'applications provenant d'autres domaines ne peut se faire,
en particulier dans les domaines physiques, où les systèmes d'équations
deviennent de plus en plus complexes.

Pendant ma thèse, j'ai développé une méthode algorithmique itérative pour
résoudre des approximations de rang faible sur des matrices. Afin de rendre
cette méthode efficace en temps de calcul, j'ai défini l'algorithme de manière
à ce qu'il s'exécute avec plusieurs précisions, le rendant ainsi applicable en
précision mixte et sur diverses architectures telles que les GPU. Étant donné
que la faible précision est généralement plus rapide que la haute précision,
j'ai obtenu une nette accélération des calculs sur différentes architectures,
notamment les CPU et les GPU.

J'ai ensuite cherché à étendre cette approche aux tenseurs. Pour cela, j'ai
d'abord implémenté un algorithme stable d'approximation de rang faible pour les
tenseurs. Une fois cette étape réalisée, j'ai combiné ces deux contributions
afin d'améliorer la gestion mémoire des tenseurs via leurs approximations de
rang faible, tout en maintenant un coût de calcul très réduit.

Dans la même thématique, mon travail de post-doctorant m'a permis de prendre
conscience qu'un problème d'optimisation issu d'une application physique réelle
ouvre de nombreuses perspectives de recherche. En effet, j'ai exploité la
structure de la matrice donnée afin d'en extraire un algorithme parallèle très
efficace, mobilisant des méthodes de parallélisation que je n'avais pas
abordées auparavant.

Je me suis ainsi détaché de l'étude du problème dans son contexte physique pour
l'aborder davantage sous l'angle de la théorie des graphes, notamment à travers
la parallélisation basée sur les tâches. Cette approche m'a permis de mettre en
évidence une stratégie de parallélisation performante, en combinant plusieurs
méthodes et en les implémentant à partir d'interfaces existantes, tout en
tirant parti de leurs avantages et en tenant compte de leurs contraintes.

\subsection{Thèse}

\begin{itemize}
    \item[\ding{118}~ Titre :] Mixed precision algorithms for low-rank matrix and tensor approximations
          \begin{itemize}
              \item[\ding{226}~ Directeur de thèse :] Marc Baboulin
                    (\href{mailto:baboulin@lmf.cnrs.fr}{baboulin@lmf.cnrs.fr})
              \item[\ding{226}~ Co-directeur de thèse :] Oguz Kaya
                    (\href{mailto:oguz.kaya@universite-paris-saclay.fr}{oguz.kaya@universite-paris-saclay.fr}),\\
                    Théo Mary (\href{mailto:theo.mary@lip6.fr}{theo.mary@lip6.fr})
          \end{itemize}
\end{itemize}

\paragraph*{Résumé}\mbox{}\\

Ma thèse propose une étude sur l'approximation de rang faible des matrices et
des tenseurs. Plus précisément nous avons défini de nouveau algorithmes stables
pour la décomposition de rang faible pour les tenseurs. Et des méthodes de
précisions mixtes pour accélérer les algorithmes en calculant en faibles
précisions tout en atteignant une haute précision d'approximation. Cette
approche sert à résoudre les problèmes d'espaces apparaissant dans les
tenseurs.

La gestion des données est souvent réalisée par des objets mathématiques tels
que les matrices et les tenseurs, qui sont la généralisation des matrices à
plus de deux dimensions. Certains domaines d'application nécessitent de stocker
trop d'éléments, créant des tenseurs trop grands ; ce problème est connu sous
le nom de \emph{curse of dimensionality}. Des méthodes mathématiques telles que
les approximations de rang faible ont été développées pour réduire la
dimensionnalité de ces objets malgré un coût très élevé en temps de calcul. De
plus, de nouvelles architectures informatiques telles que les GPU nous
permettent d'effectuer des calculs rapidement, notamment lors de calculs en
faible précision. Combiner ces nouvelles architectures avec l'approximation de
rang faible est une solution malgré la qualité des résultats altérée par la
faible précision. Ma thèse vise à proposer des algorithmes d'approximation de
rang faible stables en faible précision tout en conservant l'accélération
inhérente au calcul en faible précision, ce qui est réalisable grâce au calcul
en précision mixte.

Dans un premier temps, nous avons développé un algorithme d'approximation de
rang faible pour les tenseurs basé sur CP-ALS mais en précision mixte. En effet
sachant que la méthode CP-ALS est une méthode dite \textit{itérative}, cela
implique qu'à chaque itération nous avons une approximation qui doit être
meilleure que la précédente. Malheureusement, la non-convergence de la méthode
rendait très compliqué l'analyse de l'impact de la précision mixte sur le
résultat, on pouvait se demander si le résultat n'était pas plutôt dû à la seed
ou bien au choix du rang.

Alors pour s'assurer du bon fonctionnement de l'algorithme, j'ai pris un autre
angle d'attaque en me concentrant plutôt sur les méthodes dites
\textit{directes}. Donnant une approximation quasi-optimal et sans itérer,
cette approche avec la faible précision m'a donné des résultats très concluants
tel que montrés ci-dessous~\cref{fig.irlra}, et nous permet ainsi de définir la
méthode comme \textit{raffinement itérative} d'approximation de rang faibles.
Ce raffinement itératif s'exprime de la façon suivante,
\begin{enumerate}
    \item Calculer $\E = \X - \Decompress(\F)$ en haute précision
          $\ueps$.\label{line.decompress}
    \item If $\norm{\E} \le \eps\norm{\X}$, exit.\label{line.stop}
    \item Calculer $\FE = \LRA(\E, \epsl)$ en faible précision $\ul$.\label{line.FE}
    \item Calculer $\F = \Recompress(\F + \FE, \epsl^{i+1})$ en haute précision
          $\ueps$.\label{line.recompress}
\end{enumerate}
et atteint la borne suivante
\begin{equation}\label{eq.Fhat}
    \norm{X-\Fhat} \le (\phi^{k+1} + \xi + O(\epsl\eps))\norm{X},
\end{equation}
après $k$ itérations avec $\phi = (1+b_1\tl)\epsl + (b_2+1)\theta\eps$ la
vitesse de convergence et $\xi = \big(1+(b_2+b_3+2)\theta\big)\eps$ la
précision atteignable (prouvé dans~\cite[Théoreme 2.1]{bkmr23}).

La~\cref{graph.randSVD} représente la méthode sur une matrice Poisson (à
gauche) et un tenseur d'une application (à droite). L'abscisse représente le
nombre de \textit{raffinement itératif} qu'on fait pour atteindre la précision
souhaitée (ici double-FP64) par rapport à l'erreur relative calculée en axe des
ordonnées ainsi $\norm{A - \tilde{A}}/\norm{A}$. Chaque courbe représente la
précision dans laquelle sont faits les calculs en faibles précision et le
nombre indiqué à côté de chaque marqueur représente le rang obtenu. On constate
bien la convergence de la méthode et ce peu importe la faible précision
choisie, on constate aussi que le rang obtenu est équivalent à celui de la
précision haute ce qui implique que la méthode comprime tout aussi bien qu'une
approximation en haute précision.

Cette bonne compression n'est pas trivialle car si on ne considère pas
l'étape~\cref{line.recompress}, le rang peut être surestimé par un multiple de
3. En effet, les calculs d'algèbre linéaire montrent que le rang sans
recompression serait borné par
\begin{align*}
    \rank(\F+\F_\E)
     & \le \rank(\F) + \rank(\F_\E)                                    \\
     & = \rank(\F) + \rank(E,\epsl)                                    \\
     & \le \rank(\X,\epsl^{i+1}/2) + 2\rank(\X,\epsl^i) \label{eq.rkF} \\
     & \le 3\rank(\X,\eps).
\end{align*}

\begin{figure}[t!]
    \centering
    \begin{subfigure}{0.49\textwidth}
        \includegraphics[width=\linewidth]{fig/RandSVD_Poisson_7_10_scale_4.pdf}
        \caption{randSVD}
        \label{graph.randSVD}
    \end{subfigure}
    \begin{subfigure}{0.49\textwidth}
        \includegraphics[width=\linewidth]{fig/Otensor_ttsvd_nd6_scale_2.pdf}
        \caption{TT-SVD}
        \label{graph.ttsvd}
    \end{subfigure}
    \caption{Résultat du \textit{raffinement itératif} sur une matrice Poisson avec RandSVD comme méthode d'approximation de range faible (à gauche);  et un tenseur d'application avec TT-SVD comme méthode (à droite).}
    \label{fig.irlra}
\end{figure}

Ainsi on a taclé le problème de la \textit{curse of dimensionality} dans cette
première contribution. Seulement le calcul d'approximation de rang faible est
toujours très coûteux en temps surtout avec des matrices de grandes taille.
Alors pour pallier à ce problème, j'ai appliqué les principes de méthodes
randomisées avec toujours de la précision mixte pour accélérer les calculs et
précisemment sur des architectures GPUs. La fusion de la randomisation avec la
précision est applicable comme l'a démontré Ootomo et Yokota~cite{ooyo23}, mais
seulement sur la matrice aléatoire. J'ai alors étendu l'analyse en appliquant
la faible précision sur la multilplication entre la matrice originale et
aléatoire et exploité des nouvelles unités de calculs, nommées \textit{tensor
    cores}, pour faire des calculs en faible précision mais accumulées en haute
précision. J'ai étudié différents scénarios de ce calcul comme le
montre~\cref{graph.gemmGPU} puis j'ai constaté que le \textit{bottleneck}
devenait la méthode d'approximation.

Mais vouloir appliquer la faible précision sur la méthode d'approximation n'est
pas pertinente car cela dégraderait trop le résultat, alors j'ai changé de
noyau d'approximation en prenant une factorisation de Cholesky et ai conservé
la haute précision car la différence de temps est négligeable entre les deux
comme constaté dans~\cref{graph.cholQR}.

\begin{figure}[t!]
    \centering
    \begin{subfigure}{0.49\textwidth}
        \includegraphics[width=\linewidth]{fig/gemm_perf.pdf}
        \caption{Performance de $\Omega * A$ en précision mixte.}
        \label{graph.gemmGPU}
    \end{subfigure}
    \begin{subfigure}{0.49\textwidth}
        \includegraphics[width=\linewidth]{fig/cholQR_perf.pdf}
        \caption{Performance HouseholderQR versus CholeskyQR.}
        \label{graph.cholQR}
    \end{subfigure}
    % \caption{Résultat du \textit{raffinement itératif} sur une matrice Poisson avec RandSVD comme méthode d'approximation de range faible (à gauche);  et un tenseur d'application avec TT-SVD comme méthode (à droite).}
    \label{fig.optiGPU}
\end{figure}

En combinant ces résultats avec le raffinement itératif, nous obtenons des
résultats confortant l'apport de gain de la précision mixte sur des
architectures récentes.

\begin{figure}[t!]
    \centering
    \begin{subfigure}{0.49\textwidth}
        \includegraphics[width=\linewidth]{fig/randSVD_perf_with_raf.pdf}
        \caption{Performance.}
        \label{graph.perfirGPU}
    \end{subfigure}
    \begin{subfigure}{0.49\textwidth}
        \includegraphics[width=\linewidth]{fig/randSVD_accuracy_with_raf.pdf}
        \caption{Accuracy.}
        \label{graph.accuracyirGPU}
    \end{subfigure}
    % \caption{Résultat du \textit{raffinement itératif} sur une matrice Poisson avec RandSVD comme méthode d'approximation de range faible (à gauche);  et un tenseur d'application avec TT-SVD comme méthode (à droite).}
    \label{fig.irGPU}
\end{figure}

Pour faire la même contribution sur les tenseurs, il faut s'assurer que les
algorithmes de recompressions sont stables et cela indépendamment du format
traité. C'est pourquoi, j'ai fait une analyse de stabilité des opérations de
réseaux de tenseurs nécessaire pour le raffinement itératif. Et à partir
d'opérations stables j'ai défini une méthode de recompression stable,
démontrées par l'analyse, bien que la composition d'opérations stables
n'implique pas nécessairement un résultat stable~\cite{bnv21}. La comparaison
avec les méthodes similaires, utilisant les gramians, montrent que nous sommes
plus stables et que nous comprimons mieux qu'eux sur de faible précision.

% Nous avons développé une méthode générale d'approximation de tenseurs en
% précision mixte en calculant d'abord une approximation en faible précision et
% en l'affinant itératifment avec une précision supérieure pour maintenir la
% qualité du résultat. Sachant que cette accélération provient principalement des
% architectures GPU, plus précisément d'unités de calcul spécialisées appelées
% \emph{tensor cores}, nous avons développé une méthode générale d'approximation
% matricielle pour les architectures GPU en précision mixte utilisant ces
% \emph{tensor cores}. Notre méthode maintient la qualité du résultat, mais au
% prix d'une approximation de dimension supérieur à celle des applications
% standards. Pour compenser cet écart, des méthodes de recompression de dimension
% existent pour différents formats de tenseurs. Notre contribution finale propose
% un cadre pour l'analyse de stabilité des opérations de réseaux tenseurs
% communs, qui nous guide également vers une méthode de recompression stable.

\paragraph*{Mots-clés :} Calcul en précision mixte, Calcul haute performance, Calcul tensoriel, Réduction de données.

\subsection{Publications}

\paragraph*{Conférences internationales}\mbox{}\\[5pt]
Marc Baboulin, Simplice Donfack, Oguz Kaya, Théo Mary, and \underline{\bfseries Matthieu Robeyns}.
Mixed precision randomized low-rank approximation with GPU tensor cores.
In \textit{Mixed precision randomized low-rank approximation with GPU tensor cores}.
Proceedings of the Euro-PAR 2024 Conference, Lecture Notes in Computer Science, Springer, Vol. 14803, pp. 31-44 (2024),
\href{https://hal.science/LISN/hal-04520893v1}{Link paper.}
\mbox{}\\[10pt]

\noindent
Marc Baboulin, Oguz Kaya, Théo Mary, and \underline{\bfseries Matthieu Robeyns}.
In \textit{Mixed precision iterative refinement for low-rank matrix and tensor approximations.}
Lecture Notes in SIAM Journal on Scientific Computing, Vol. 47, Iss. 5 (2025), pp. A2403-C1171
\href{https://hal.science/hal-04115337v1/file/Paper.pdf}{Link paper}.
\mbox{}\\[10pt]

\noindent
Marc Baboulin, Oguz Kaya, Théo Mary, and \underline{\bfseries Matthieu Robeyns}.
In \textit{Numerical stability of tree tensor network operations, and a stable rounding algorithm.}
Lecture Notes in SIAM Journal on SIMAX [Under review]
\href{https://hal.science/hal-04996127}{Link paper}.
\mbox{}\\[10pt]

\paragraph*{Présentations internationales}\mbox{}\\[5pt]
Marc Baboulin, Oguz Kaya, Théo Mary, and \underline{\bfseries Matthieu Robeyns}.
In \textit{SIAM CSE, Amsterdam, Feburary 2023} - Mixed Precision Iterative Refinement for Low-Rank Matrix and Tensor Approximations.
\mbox{}\\[10pt]

\noindent
Marc Baboulin, Oguz Kaya, Théo Mary, and \underline{\bfseries Matthieu Robeyns}.
In \textit{Euro-PAR, Madrid, August 2024} - Mixed precision randomized low-rank approximation with GPU tensor cores.
\mbox{}\\[10pt]

\noindent
Marc Baboulin, Oguz Kaya, Théo Mary, and \underline{\bfseries Matthieu Robeyns}.
In \textit{HPCSIM, Padova, September 2025} - A parallel approach for solving the Green's function based on the nested-dissection a and selected-inverse methods.
\mbox{}\\[10pt]

\paragraph*{Présentation nationale}\mbox{}\\[5pt]
Marc Baboulin, Oguz Kaya, Théo Mary, and \underline{\bfseries Matthieu Robeyns}.
In \textit{SIAM LA, Paris, May 2024} - A general framework for the stable rounding of low-rank tensors with error analysis.
\mbox{}\\[10pt]

\subsection{Projet EoCoE-III (Post-doctorat)}

Je suis, depuis Janvier 2025, Post-Doctorant à l'ENSEEITH Toulouse, dans
l'équipe APO. Je travaille sur le projet
\href{https://www.eurohpc-ju.europa.eu/research-innovation/our-projects/eocoe-iii_en}{EuroHPC
    EoCoE-III}, mené par trois pays d'Europe : la France, l'Allemagne et l'Italie.
Ce projet vise à utiliser les supercalculateurs européens pour les analyses de
transports de particules, permettant ainsi la génération de nouveaux circuits
quantiques. Les nouveaux circuits étant trop coûteux en ressources mémoire et
en temps, mon objectif au sein du projet est d'optimiser les algorithmes et les
calculs mathématiques via des architectures très performantes, récentes.

% Alors pour extraire du parallélime, je dois me ramener à une approche à base de
% graph et penser à permuter les blocks de façon efficace. Cette idée provient de
% la \textit{nested dissection} qui me permet ainsi de parallélisé le problème à
% un certains niveau souhaité. En contrepartie de cette parallélisation,
% j'augmente le nombre de calculs par un effet de \textit{fill-in} que je peux
% contrôler sachant que la matrice étudié est tridiagonal.

% L'optimisation est aussi présente au sein des calculs
% employé dans les algorithmes et de comment rendre plus abordables des calculs
% d'inverses de matrices sur ordinateurs (via l'algèbre linéaire). La
% contribution va mener à un papier prévu pour cette année.

\subsubsection*{Résumé}

L'application vise à résoudre la~\gls{negf}~\cite{sch61} équation possédant
deux inconnus $G^r$ et $G^n$. $G^r$ provient de la Dyson équation et peut
synthétiquement s'écrire comme $T G^r = I$ où $T$ est une matrice tridiagonale
par bloc et $I$ la matrice identité. Sachant que calculer $G^r = T^{-1}$
produit une matrice dense (ce qui n'est pas convenable), une solution fût de
calculer une approximation de $T^{-1}$ en tirant avantages de la structure
tridiagonal de la matrice $T$.

Un moyen de rendre le problème moins complexe est de commencer par le
factoriser, puis de calculer l'inverse de ces facteurs. J'optimise la méthode
en profitant de la structure tridiagonale du problème afin de ramener le calcul
de l'inverse à une \textit{selected inverse}~\cite{lyml11}, consistant à ne
calculer que les blocs non nuls de l'inverse en utilisant uniquement les blocs
non nuls. Ainsi, on évite des calculs inutiles ainsi que des surcoûts en
mémoire. Comme dans ma thèse, le solveur reste toutefois coûteux en temps de
calcul et purement séquentiel en raison de la dépendance entre les blocs
diagonaux.

% L'algorithme utilisé permettant ceci est la~\gls{si}~\cite{lyml11} qui
% n'utilise uniquement les blocs non nuls pour calculer uniquement l'inverse des
% blocs de la matrice $T$; qu'on factorise au préalable pour une meilleure
% efficacité. Bien que l'algorithme soit efficace, celui-ci reste limité par son
% exécution purement séquentielle.

Alors, la seconde contribution fût de paralléliser la méthode en ajoutant une
couche d'abstraction par la méthode de dissection emboitée
(~\gls{nd}\cite{george73}). Cette approche consiste à subdiviser la matrice en
plusieurs domaines de calcul au moyen d'une permutation de blocs spécifiques,
appelés séparateurs (e.g. \textit{separators}). Ainsi, une bonne scalabilité
peut être atteinte, malgré un surcoût de calcul induit par la permutation de la
matrice ; ce phénomène est appelé l'effet de remplissage (e.g.,
\textit{fill-in} effect). Afin de minimiser le taux de remplissage, une autre
stratégie de parallélisation a été utilisée, elle consiste à diviser les blocs
de la matrice en plus petits blocs, nommés tuiles (e.g. \textit{tiles}) donnant
accès à une plus petite granularité des calculs. Cette méthode nommée
\textit{block splitting}~\cite{pabo14} peut aussi être assemblée avec
la~\gls{nd} pour maximiser la charge de calcul en fonction de l'architecture
utilisée. La~\cref{graph.methodPostDoc} illustre les différentes méthodes
employées pour calculer l'inverse de la matrice $T$, présentées de haut en bas.

Afin d'exploiter efficacement ces deux méthodes de parallélisation,
l'utilisation conjointe des modèles \textit{task-based} et \textit{sequential
    task flow} s'avère particulièrement adaptée. Elle permet de générer du code
parallèle avec peu de modifications du code séquentiel, tout en assurant un
contrôle explicite des dépendances entre les tâches. L'exécution des
algorithmes peut désormais être représentée sous forme d'arbres binaires,
parcourus dans un sens pour la factorisation et dans l'autre pour le calcul de
l'inverse. La~\cref{fig.daglu} illustre l'arbre de factorisation pour deux
niveaux de dissection emboîtée, où chaque couleur correspond à un noyau de
calcul utilisé dans l'algorithme. Il convient de noter que le chemin critique
de cet arbre correspond toujours à l'une des branches les plus internes.

\begin{figure*}
    \centering
    \begin{subfigure}{0.49\textwidth}
        \includegraphics[width=\linewidth]{fig/postDocResume.png}
        \caption{Différentes méthodes appliqué sur la matrice d'étude $A$ (à gauche) pour trouver son inverse (à droite).}
        \label{graph.methodPostDoc}
    \end{subfigure}
    \begin{subfigure}{0.49\textwidth}
        \includegraphics[width=\linewidth]{fig/Daglu.png}
        \caption{Directed acyclic graph (DAG) de la factorisation LU de la matrice $T$.}
        \label{fig.daglu}
    \end{subfigure}
    \caption{Modélisation du calcul de l'inverse (à gauche) et du graphe des tâches (à droite).}
\end{figure*}

L'implémentation fût initialement faite en \textbf{Julia}, mais un souci de
surcoût dans la librairie
\textbf{Dagger}~\href{https://juliaparallel.org/Dagger.jl/stable/}{$^{[dagger]}$}
a fait basculer l'implémentation dans le langage \textbf{C} avec les librairies
\textbf{Chameleon}~\href{https://solverstack.gitlabpages.inria.fr/chameleon/}{$^{[cham]}$}
et \textbf{StarPU}\href{https://starpu.gitlabpages.inria.fr/}{$^{[starpu]}$}
pour le solver et l'ordonnancement, respectivement. Une intégration avec la
bibliothèque
\textbf{libNEGF}~\href{https://github.com/libnegf/libnegf}{$^{[negf]}$} codée en
\textbf{Fortran} est prévue d'ici la fin de l'année.

\begin{figure}[t!]
    \centering
    \begin{subfigure}{0.49\textwidth}
        \includegraphics[width=\linewidth]{fig/plot_heat_dgetri_nd_time.pdf}
        \caption{Nested dissection.}
        \label{fig.perfND}
    \end{subfigure}
    \begin{subfigure}{0.49\textwidth}
        \includegraphics[width=\linewidth]{fig/plot_time_subcut.pdf}
        \caption{Nested dissection + Block splitting.}
        \label{fig.perfNDBS}
    \end{subfigure}
    \caption{Performances des deux méthodes étudiées pendant le projet.}
    \label{fig.perfPD}
\end{figure}

Les epxériences faites sur un cluster de l'Inria Bordeaux nommé
\textbf{PlaFRIM}~\href{https://www.plafrim.fr/}{$^{[plafrim]}$} nous ont permis
de générer des résultats sur les performances de la nested
dissection~\cref{fig.perfND} et de la nested dissection combinée avec le block
splitting~\cref{fig.perfNDBS}.
% Les résultats de ce projet sont pour des architectures homogènes ou hétérogènes
% (CPU/GPU) et utilisent majoritairement des opérations de type
% BLAS-3~\href{https://www.netlib.org/blas/}{$^{[blas]}$} optimal pour ces
% architectures et on était testé sur un cluster de l'Inria Bordeaux nommé
% \textbf{PlaFRIM}~\href{https://www.plafrim.fr/}{$^{[plafrim]}$}.

\printbibliography[keyword=recherche]